\subsection{Velocity, central vorticity and the maximal time}

\begin{proposition}\label{cc:uf:terminal:velocity-curl}
The solution has bounded velocity on $[0,T_\infty)$, and
$|\curl u(t,0)|\longrightarrow\infty$ as $t\uparrow T_\infty$.
\end{proposition}
\begin{proof}
The continuous field $w$ on $[0,T_\infty]\times\mathbb T^3$ is bounded,
and Proposition~\ref{cc:uf:terminal:agreement} identifies it with $u$
before $T_\infty$.

For the vorticity, fix a large $j$ and consider the whole interval
$[t_j,t_{j+1}]$. The preceding wave is fully active there. Applying
\textup{(T4)} of Theorem~\ref{cc:uf:tower:thm:finite} at steps $j-1$
and $j$ gives, at the origin,
\[
 \nabla U_j=\nabla U_{j-2}
   +\mathbf a_{j-1}\otimes\kappa_{j-1}
   +\widetilde\chi_j\,\mathbf a_j\otimes\kappa_j+E_{c,j-1}+E_{c,j}.
\]
A gradient $\mathbf a\otimes\kappa$ contributes $\kappa\times\mathbf a$
to vorticity. Projecting $\curl U_j$ along $\kappa_j$ therefore removes
the leading curl of wave $j$. The preceding wave contributes
$\kappa_{j-1}\times\mathbf a_{j-1}=\sigma_{j-1}E_3$ in its own frame.
By \eqref{cc:eq:scaled-pair} and Lemma~\ref{cc:uf:exact:lem:additive},
the $E_3$ component of $\kappa_j/\lambda_j$ is at least $1/2$ and its
length is at most $C(1+\tau^2)$ at normalized time $\tau$.
The normalized length of $[t_j,S_j]$ is at most $CL_{j+1}^4$ by item 2
of the proof of Theorem~\ref{cc:uf:tower:thm:finite}. Thus
\begin{equation}\label{cc:uf:terminal:eq:projection}
 \frac{\kappa_j}{|\kappa_j|}\cdot(\kappa_{j-1}\times\mathbf a_{j-1})
    \ge c\sigma_{j-1}L_{j+1}^{-8},\qquad
 \frac{\kappa_j}{|\kappa_j|}\cdot(\kappa_j\times\mathbf a_j)=0.
\end{equation}

The preceding wave retains its strength throughout $[t_j,t_{j+1}]$.
Put $\mu=\lambda_{j-1}$. At $t_j$, \eqref{cc:uf:tower:eq:late-values}
gives $\sigma_{j-1}\ge c\mu^{1/4}$. By \eqref{cc:eq:frame},
$|\partial_t\log\sigma_{j-1}|\le|\nabla U_{j-2}(t,0)|
\le\lambda_{j-2}^9\le\mu^{1/32}$, using
\eqref{cc:uf:tower:eq:derived-low} and item 2 above. Since the interval
has length at most $\tau_j\le\mu^{-1/16}$ by \textup{(T4)} and
\eqref{cc:uf:tower:eq:increment-tail}, $\log\sigma_{j-1}$ changes by
at most one. Hence $\sigma_{j-1}\ge c\mu^{1/4}$ on the whole interval.
Since $L_{j+1}^8\le(\log\mu)^{16}\le\mu^{1/8}$ for large $j$ by
\eqref{cc:uf:tower:eq:uniform}, the projected curl of the preceding wave
in \eqref{cc:uf:terminal:eq:projection} is at least $c\mu^{1/8}$.
The older gradient contributes at most $C\lambda_{j-2}^9$ to the curl
and the two central errors are at most $2\mu^{-1/4}$ in norm. Consequently
\begin{equation}\label{cc:uf:terminal:eq:curl-lower}
 |\curl U_j(t,0)|
 \ge c\lambda_{j-1}^{1/8}-C\lambda_{j-2}^9-C,
\end{equation}
which tends to infinity uniformly on $[t_j,t_{j+1}]$.

For $\ell>j$ the oscillation $O_\ell$ vanishes throughout $[t_j,t_{j+1}]$.
Applying \eqref{cc:uf:tower:eq:increment} to the remaining increments
at order one and taking the limit by Proposition~\ref{cc:uf:terminal:agreement}
gives
\[
 \sup_{t_j\le t\le t_{j+1}}|\curl(u-U_j)(t,0)|
 \le C\sum_{\ell>j}\lambda_\ell^{-1}.
\]
This tail tends to zero. The successive intervals cover
$[t_j,T_\infty)$, so \eqref{cc:uf:terminal:eq:curl-lower} gives
$|\curl u(t,0)|\to\infty$ as $t\uparrow T_\infty$.
\end{proof}

\begin{corollary}\label{cc:uf:terminal:time}
The maximal existence time is $T_*=T_\infty\in(0,\infty)$, and
both integrals in \eqref{cc:uf:eq:main-breakdown} diverge.
\end{corollary}
\begin{proof}
By Proposition~\ref{cc:uf:terminal:agreement}, $T_*\ge T_\infty$. If
$T_*>T_\infty$, then smoothness on the compact interval $[0,T_\infty]$
would bound $\curl u(t,0)$, contrary to
Proposition~\ref{cc:uf:terminal:velocity-curl}. Hence $T_*=T_\infty$,
which is positive and finite by Theorem~\ref{cc:uf:tower:thm:finite}.
The initial data are smooth, so the periodic Beale--Kato--Majda criterion
\cite{BKM}, in the formulation of \cite[Theorem~1]{GGPP}, implies
$\int_0^{T_*}\|\curl u(t)\|_\infty\dd t=\infty$.
The gradient integral also diverges, since
$\|\curl u\|_\infty\le C\|\nabla u\|_\infty$.
\end{proof}

\subsection{The velocity away from the origin at the singular time}

\begin{proposition}\label{cc:uf:terminal:off-origin}
For every compact $K\subset\mathbb T^3\setminus\{0\}$ and every integer
$m\ge0$, both $u(t)$ and $\partial_tu(t)$ converge in $C^m(K)$ as
$t\uparrow T_\infty$.
\end{proposition}
\begin{proof}
Fix $K$ and $m$ as in the statement. The support of $O_j$ lies in the
ball of radius $\lambda_{j-1}^{-1}$ about the origin throughout $[0,T_\infty]$.
For large $j$ the ball is disjoint from $K$ and $q_j\ge m$, so
\eqref{cc:uf:tower:eq:increment} gives
$\|U_j-U_{j-1}\|_{C([0,T_\infty];C^m(K))}\le\lambda_j^{-1}$.
The increments are summable in this norm, hence $u(t)\to w(T_\infty)$
in $C^m(K)$.

The time derivative also contains the pressure gradient, which depends
on the velocity on the whole torus. Choose a smooth cutoff $\chi$ equal
to one near $K$ and zero near the origin. On $[0,T_\infty)$ we have
\[
 \partial_tu=-\PP\diver(\chi u\otimes u)-\PP\diver\bigl((1-\chi)u\otimes u\bigr).
\]
Applying the spatial convergence just proved on $\supp\chi$ through
order $m+2$ gives convergence of $\chi u\otimes u$ in $\Holder{m+1}$
on the torus. Since $\PP\diver$ maps $\Holder{m+1}$ boundedly to
$\Holder{m}$, the first term converges in $C^m(K)$.

For the second term, $\supp(1-\chi)$ is a positive distance from $K$.
Only the pressure part $-\nabla\Delta^{-1}\diver\diver$ acts across
this distance. Its kernel is smooth off the diagonal, so for any
bounded tensor field $A$,
\[
 \|\PP\diver((1-\chi)A)\|_{C^m(K)}\le C_{m,K,\chi}\|A\|_{L^1}.
\]
Since $u(t)\otimes u(t)\to w(T_\infty)\otimes w(T_\infty)$ uniformly
on the torus, the difference tends to zero in $L^1$. The estimate
therefore gives convergence of the second term in $C^m(K)$.
\end{proof}

Propositions~\ref{cc:uf:terminal:agreement},
\ref{cc:uf:terminal:velocity-curl} and~\ref{cc:uf:terminal:off-origin}
and Corollary~\ref{cc:uf:terminal:time} prove
Theorem~\ref{cc:target:unforced}.
